\documentstyle{article}

\title{The Parameterization Problem for of Algebraic Surfaces}
\author{Josef Schicho%
        \thanks{supported by the Austrian Science Fund in the
                frame of project P12662-TEC} \\
        Research Institute for Symbolic Computation \\
        Johannes Kepler University \\
        A-4040 Linz, Austria \\
        e-mail: {\tt Josef.Schicho@risc.uni-linz.ac.at} \\
        www: {\tt http://www.risc.uni-linz.ac.at/people/jschicho}}
\newtheorem{thm}{Theorem}
\newtheorem{prop}{Proposition}
\newtheorem{cor}{Corollary}
\newtheorem{lemma}{Lemma}
\newtheorem{rem}{Remark}
\newtheorem{example}{Example}
\newtheorem{definition}{Definition}
%\hyphenation{pa-ra-met-ri-zation pa-ra-met-ric}
\begin{document}
\maketitle

\section{Extended Abstract}

We give an overview on the parametrization problem, including recent
achievements and open problems.

The parameterization problem is the following. Given is the implicit
representation of a surface, in terms of an algebraic equation
in three variables.
Wanted are three rational functions in two variables,
such that the image of the rational map defined by
this triple is a two-dimensional subset of the surface.
For instance, the the unit sphere
with equation $x^2+y^2+z^2-1=0$
has the parameterization
\[ (x,y,z) =
  (\frac{2s}{s^2+t^2+1},\frac{2t}{s^2+t^2+1},\frac{s^2+t^2-1}{s^2+t^2+1}) , \]
which can be found by stereographic projection.

Over the complex numbers, the well-known criterion of
Castelnouvo~\cite{Castelnuovo:96} is very relevant: a complex
surface has a rational parametrization iff its arithmetic genus
and its second plurigenus both vanish. An algorithm for deciding
this criterion and computing a parametrization in case the criterion
is fulfilled is contained in \cite{Schicho:97}. This algorithm
is currently being implemented by W. Hei{\ss} and the author,
and we give a first report on computational experiments.

Over the real numbers, Castelnuovo's criterion is necessary,
but it is not known whether it is sufficient
for the existence of a parametrization with real coefficients.
Important subcases have been solved, for instance the case
of tubular or conical surfaces (see \cite{Schicho:98c,Peternell:97}).
This result is relevant for applications in CAD/CAM
(see \cite{Arrondro_Sendra_Sendra:97,Pottmann:97}).

Related problems are the parametrization problem for algebraic
curves (studied in
\cite{Walker:78,Abhyankar_Bajaj:87,Sendra_Winkler:91,Schicho:92,%
Hoeij:94,Mnuk_Sendra_Winkler:96,Mnuk:96,Sendra_Winkler:97,Hoeij:97})
and the implicitization problem
studied in \cite{Canny_Manocha:92,Gao_Chou:92,Kalkbrener:90c,Choinh_Goldman:92}.
Another related problem is the problem of resolution of singularities
(see \cite{Walker:35,Hironaka:64,Lipman:78,Villamayor:89,Bierstone_Milman:91,%
Bodnar_Schicho:98,Schicho_u:99c}),
which is required as a subtask for surface parametrization.



\begin{thebibliography}{10}

\bibitem{Abhyankar_Bajaj:87}
{\sc Abhyankar, S.~S., and Bajaj, B.}
\newblock {Automatic parametrization of curves and surfaces III}.
\newblock {\em Computer Aided Geometric Design 5-4\/} (1987), 309--323.

\bibitem{Arrondro_Sendra_Sendra:97}
{\sc Arrondo, E., Sendra, J., and Sendra, J.~R.}
\newblock Parametric generalized offsets to hypersurfaces.
\newblock {\em Journal of Symbolic Computation 23\/} (1997), 267--285.

\bibitem{Bierstone_Milman:91}
{\sc Bierstone, E., and Milman, P.}
\newblock A simple constructive proof of canonical resolution of singularities.
\newblock In {\em Effective methods in algebraic geometry}, T.~Mora and
  C.~Traverso, Eds. {Birkh\"{a}user}, 1991, pp.~11--30.

\bibitem{Schicho_u:99c}
{\sc Bodnar, G., and Schicho, J.}
\newblock Automated resolution of singularities for hypersurfaces.
\newblock Tech. Rep. 99-06, RISC-Linz, Univ. Linz, A-4040 Linz, 1999.

\bibitem{Bodnar_Schicho:98}
{\sc Bodn\'{a}r, G., and Schicho, J.}
\newblock A computer program for the resolution of singularities.
\newblock In {\em Working Week on Mountains and Singularities\/} (1999),
  Springer.
\newblock to appear.

\bibitem{Canny_Manocha:92}
{\sc Canny, J., and Manocha, D.}
\newblock Implicit representation of rational parametric surfaces.
\newblock {\em Journal of Symbolic Computation 13}, 5 (1992).

\bibitem{Castelnuovo:96}
{\sc Castelnuovo, G.}
\newblock {Sulle superficie di genere zero}.
\newblock In {\em {Memoria scelte}}. Zanichelli, 1939, pp.~307--334.

\bibitem{Choinh_Goldman:92}
{\sc Chionh, E.-W., and Goldman, R.}
\newblock {Degree, multiplicity and inversion formulas for rational surfaces
  using $u$-resultants}.
\newblock {\em Comp. Aided Geom. Des. 9\/} (1992), 93--108.

\bibitem{Gao_Chou:92}
{\sc Gao, X.-S., and Chou, S.-C.}
\newblock {Implicitization of Rational Parametric Equations}.
\newblock {\em J. Symb. Comput. 14\/} (1992), \ pages 459--470.

\bibitem{Hironaka:64}
{\sc Hironaka, H.}
\newblock Resolution of singularities of an algebraic variety over a field of
  characteristic 0.
\newblock {\em Ann. Math. 79\/} (1964), 109--326.

\bibitem{Kalkbrener:90c}
{\sc Kalkbrener, M.}
\newblock {Implicitization of Rational Curves and Surfaces}.
\newblock In {\em Lect. Notes in Comp. Sci. 508, AAECC-8\/} (Tokyo, Japan, Aug.
  1990), Sakata, Ed.

\bibitem{Lipman:78}
{\sc Lipman, J.}
\newblock Desingularization of two-dimensional schemes.
\newblock {\em Ann. Math. 107\/} (1978), 151--207.

\bibitem{Mnuk:96}
{\sc {M\~{n}uk}, M.}
\newblock {\em {Algebraic parametrization of rational curves}}.
\newblock PhD thesis, RISC Linz, 1996.

\bibitem{Mnuk_Sendra_Winkler:96}
{\sc {M\~{n}uk}, M., Sendra, J.~R., and Winkler, F.}
\newblock {On the complexity of parametrizing curves}.
\newblock {\em Beitr. Algebra und Geometrie 37/2\/} (1996), 309--328.

\bibitem{Peternell:97}
{\sc Peternell, M.}
\newblock {\em Rational parametrizations for envelopes of quadric families}.
\newblock PhD thesis, Techn. Univ. Vienna, 1997.

\bibitem{Pottmann:97}
{\sc Peternell, M., and Pottmann, H.}
\newblock Computing rational parametrizations of canal surfaces.
\newblock {\em Journal of Symbolic Computation 23\/} (1997), 255--266.

\bibitem{Schicho:92}
{\sc Schicho, J.}
\newblock On the choice of pencils in the parametrization of curves.
\newblock {\em J. Symb. Comp. 14}, 6 (1992), 557--576.

\bibitem{Schicho:98c}
{\sc Schicho, J.}
\newblock Rational parameterization of real algebraic surfaces.
\newblock In {\em ISSAC-98\/} (1998), ACM Press, pp.~302--308.

\bibitem{Schicho:97}
{\sc Schicho, J.}
\newblock Rational parametrization of surfaces.
\newblock {\em J. Symb. Comp. 26}, 1 (July 1998), 1--30.

\bibitem{Sendra_Winkler:91}
{\sc Sendra, J., and Winkler, F.}
\newblock Symbolic parametrization of curves.
\newblock {\em Journal of Symbolic Computation 12}, 6 (1991), 607--632.

\bibitem{Sendra_Winkler:97}
{\sc Sendra, J., and Winkler, F.}
\newblock Parametrization of algebraic curves over optimal field extensions.
\newblock {\em Journal of Symbolic Computation 23}, 2/3 (1997), 191--208.

\bibitem{Hoeij:94}
{\sc van Hoeij, M.}
\newblock {Computing parametrizations of rational algebraic curves}.
\newblock In {\em ISSAC-94\/} (1994), ACM Press, pp.~187--190.

\bibitem{Hoeij:97}
{\sc van Hoeij, M.}
\newblock {Rational parametrizations of algebraic curves using canonical
  divisors}.
\newblock {\em Journal of Symbolic Computation 23\/} (1997), 209--227.

\bibitem{Villamayor:89}
{\sc Villamayor, O.}
\newblock {Constructiveness of Hironaka's resolution}.
\newblock {\em Ann. Scient. Ecole Norm. Sup. 4 22\/} (1989), 1--32.

\bibitem{Walker:35}
{\sc Walker, R.~J.}
\newblock {Reduction of the singularities of an algebraic surface}.
\newblock {\em Ann. Math. 36\/} (1935), 336--365.

\bibitem{Walker:78}
{\sc Walker, R.~J.}
\newblock {\em Algebraic Curves}.
\newblock Springer, 1978.

\end{thebibliography}

\end{document}
